\section{Übung 8 – Bedingte Wahrscheinlichkeiten und Satz von Bayes}

In dieser Übung werden bedingte Wahrscheinlichkeiten beim Ziehen ohne Zurücklegen
(Teil~A) sowie die Binomialverteilung beim Ziehen mit Zurücklegen (Teil~B)
betrachtet.

% ----------------------------------------------------------------
\subsection*{Teil A – Urne mit farbigen Kugeln (ohne Zurücklegen)}
% ----------------------------------------------------------------

Eine Urne enthält $2$ rote, $5$ schwarze und $3$ weiße Kugeln, insgesamt also
$N = 10$ Kugeln. Es wird zweimal \emph{ohne Zurücklegen} gezogen.

% ---- A1 --------------------------------------------------------
\subsubsection*{A1 – Bedingte Wahrscheinlichkeiten für die 2.\ Ziehung}

\begin{beispiel}[Bedingte Wahrscheinlichkeit nach einer schwarzen Kugel]

\textbf{Gegeben:} Urne mit $2$ roten, $5$ schwarzen, $3$ weißen Kugeln
($N = 10$). Die erste gezogene Kugel ist schwarz und wird \emph{nicht}
zurückgelegt.

\textbf{Gesucht:} Die bedingten Wahrscheinlichkeiten
$P(\text{rot} \mid S_1)$, $P(\text{schwarz} \mid S_1)$ und
$P(\text{weiß} \mid S_1)$ für die zweite Ziehung, wobei $S_1$ das Ereignis
„erste Kugel schwarz`` bezeichnet.

\textbf{Formel} (bedingte Wahrscheinlichkeit als Laplace-Experiment auf der
reduzierten Urne, Abschnitt~\ref{sec:bedingte-wahrscheinlichkeit}):

$$
P(\text{Farbe} \mid S_1) = \frac{\text{Anzahl Kugeln dieser Farbe}}{\text{Restzahl Kugeln}}
$$

\textbf{Rechnung:}

Nach dem Entnehmen einer schwarzen Kugel verbleiben $N' = 9$ Kugeln, davon
$2$ rote, $4$ schwarze und $3$ weiße. Somit gilt:

$$
P(\text{rot} \mid S_1) = \frac{2}{9}, \qquad
P(\text{schwarz} \mid S_1) = \frac{4}{9}, \qquad
P(\text{weiß} \mid S_1) = \frac{3}{9} = \frac{1}{3}
$$

Probe: $\dfrac{2}{9} + \dfrac{4}{9} + \dfrac{3}{9} = \dfrac{9}{9} = 1$ \checkmark

\textbf{Lösung:}

$$
\boxed{P(\text{rot} \mid S_1) = \frac{2}{9} \approx 0{,}222, \quad
P(\text{schwarz} \mid S_1) = \frac{4}{9} \approx 0{,}444, \quad
P(\text{weiß} \mid S_1) = \frac{1}{3} \approx 0{,}333}
$$

\end{beispiel}

% ---- A2 --------------------------------------------------------
\subsubsection*{A2 – Wahrscheinlichkeit, dass die 2.\ Kugel schwarz ist}

\begin{beispiel}[Satz von der totalen Wahrscheinlichkeit]

\textbf{Gegeben:} Urne wie oben ($2$ rot, $5$ schwarz, $3$ weiß, $N = 10$),
zweimal ohne Zurücklegen.

\textbf{Gesucht:} $P(S_2)$ -- die Wahrscheinlichkeit, dass die \emph{zweite}
gezogene Kugel schwarz ist (ohne Kenntnis der ersten Kugel).

\textbf{Formel} (Satz von der totalen Wahrscheinlichkeit,
Abschnitt~\ref{sec:totale-wahrscheinlichkeit}):

$$
P(S_2) = \sum_{i} P(S_2 \mid F_i)\, P(F_i)
$$

wobei $F_i \in \{\text{rot}, \text{schwarz}, \text{weiß}\}$ die Farbe der ersten
Kugel ist.

\textbf{Rechnung:}

Die a-priori-Wahrscheinlichkeiten der ersten Ziehung sind
$P(R_1) = \tfrac{2}{10}$, $P(S_1) = \tfrac{5}{10}$, $P(W_1) = \tfrac{3}{10}$.
Die bedingten Wahrscheinlichkeiten für eine schwarze zweite Kugel:

$$
P(S_2 \mid R_1) = \frac{5}{9}, \qquad
P(S_2 \mid S_1) = \frac{4}{9}, \qquad
P(S_2 \mid W_1) = \frac{5}{9}
$$

Einsetzen:

$$
P(S_2) = \frac{5}{9}\cdot\frac{2}{10} + \frac{4}{9}\cdot\frac{5}{10}
       + \frac{5}{9}\cdot\frac{3}{10}
       = \frac{10}{90} + \frac{20}{90} + \frac{15}{90}
       = \frac{45}{90}
$$

\textbf{Lösung:}

$$
\boxed{P(S_2) = \frac{1}{2} = 0{,}5}
$$

Die zweite Kugel ist mit derselben Wahrscheinlichkeit schwarz wie die erste
($P(S_1) = \tfrac{5}{10} = \tfrac{1}{2}$). Dieses Symmetrieergebnis gilt
allgemein: Ohne Information über die erste Ziehung ist die Wahrscheinlichkeit
für jede Position dieselbe.

\end{beispiel}

% ---- A3 --------------------------------------------------------
\subsubsection*{A3 – Rückschluss auf die erste Kugel (Satz von Bayes)}

\begin{beispiel}[Satz von Bayes]

\textbf{Gegeben:} Aus der Urne ($2$ rot, $5$ schwarz, $3$ weiß) wird eine Kugel
\emph{verdeckt} gezogen und beiseitegelegt. Anschließend wird eine schwarze
Kugel gezogen (Ereignis $S_2$).

\textbf{Gesucht:} $P(S_1 \mid S_2)$ -- die Wahrscheinlichkeit, dass die
zuerst gezogene (verdeckte) Kugel ebenfalls schwarz war.

\textbf{Formel} (Satz von Bayes, Abschnitt~\ref{sec:satz-von-bayes}):

$$
P(S_1 \mid S_2) = \frac{P(S_2 \mid S_1)\, P(S_1)}{P(S_2)}
$$

\textbf{Rechnung:}

Die benötigten Größen sind aus A1 und A2 bekannt:

$$
P(S_2 \mid S_1) = \frac{4}{9}, \qquad
P(S_1) = \frac{5}{10} = \frac{1}{2}, \qquad
P(S_2) = \frac{1}{2}
$$

Einsetzen:

$$
P(S_1 \mid S_2) = \frac{\dfrac{4}{9}\cdot\dfrac{1}{2}}{\dfrac{1}{2}}
               = \frac{4}{9}
$$

\textbf{Lösung:}

$$
\boxed{P(S_1 \mid S_2) = \frac{4}{9} \approx 0{,}444}
$$

Da $P(S_2) = P(S_1)$ ist, vereinfacht sich der Bayes-Ausdruck zu
$P(S_1 \mid S_2) = P(S_2 \mid S_1) = \tfrac{4}{9}$. Die Information, dass die
zweite Kugel schwarz ist, senkt die Wahrscheinlichkeit für eine schwarze erste
Kugel von $\tfrac{1}{2}$ auf $\tfrac{4}{9}$, da nun eine schwarze Kugel
„verbraucht`` ist.

\end{beispiel}

% ----------------------------------------------------------------
\subsection*{Teil B – Zweite Urne (mit Zurücklegen, Binomialverteilung)}
% ----------------------------------------------------------------

Eine zweite Urne enthält grüne und blaue Kugeln. Die Wahrscheinlichkeit, eine
grüne Kugel zu ziehen, beträgt $P_g$. Es wird $n = 16$ Mal eine Kugel gezogen
und jeweils wieder \emph{zurückgelegt}. Die einzelnen Ziehungen sind damit
statistisch unabhängig und identisch verteilt (Bernoulli-Kette).

% ---- B1 --------------------------------------------------------
\subsubsection*{B1 – Genau $k$-mal grün}

\begin{beispiel}[Binomialverteilung]

\textbf{Gegeben:} $n = 16$ unabhängige Ziehungen mit Zurücklegen,
Erfolgswahrscheinlichkeit $P_g$ für „grün`` pro Ziehung.

\textbf{Gesucht:} $P(X = k)$ -- die Wahrscheinlichkeit, dass genau $k$-mal eine
grüne Kugel gezogen wurde.

\textbf{Formel} (Binomialverteilung, Abschnitt~\ref{sec:binomialverteilung}):

$$
P(X = k) = \binom{n}{k}\, P_g^{\,k}\,(1 - P_g)^{\,n-k}
$$

\textbf{Rechnung:}

Mit $n = 16$ ergibt sich unmittelbar:

$$
P(X = k) = \binom{16}{k}\, P_g^{\,k}\,(1 - P_g)^{\,16-k},
\qquad k \in \{0, 1, \ldots, 16\}
$$

\textbf{Lösung:}

$$
\boxed{P(X = k) = \binom{16}{k}\, P_g^{\,k}\,(1 - P_g)^{\,16-k}}
$$

\end{beispiel}

% ---- B2 --------------------------------------------------------
\subsubsection*{B2 – Höchstens eine grüne Kugel bei $P_g = 0{,}4$}

\begin{beispiel}[Kumulierte Binomialwahrscheinlichkeit]

\textbf{Gegeben:} $n = 16$, $P_g = 0{,}4$, also $1 - P_g = 0{,}6$.

\textbf{Gesucht:} $P(X \le 1)$ -- die Wahrscheinlichkeit, dass höchstens eine
grüne Kugel gezogen wurde.

\textbf{Formel} (Summe der Einzelwahrscheinlichkeiten,
Abschnitt~\ref{sec:binomialverteilung}):

$$
P(X \le 1) = P(X = 0) + P(X = 1)
$$

\textbf{Rechnung:}

$$
P(X = 0) = \binom{16}{0}\,(0{,}4)^{0}\,(0{,}6)^{16}
         = (0{,}6)^{16} \approx 0{,}000282
$$

$$
P(X = 1) = \binom{16}{1}\,(0{,}4)^{1}\,(0{,}6)^{15}
         = 16 \cdot 0{,}4 \cdot (0{,}6)^{15} \approx 0{,}003009
$$

Summe:

$$
P(X \le 1) \approx 0{,}000282 + 0{,}003009 = 0{,}003291
$$

\textbf{Lösung:}

$$
\boxed{P(X \le 1) \approx 0{,}00329 \approx 0{,}33\,\%}
$$

Bei einer recht hohen Grün-Wahrscheinlichkeit von $40\,\%$ pro Ziehung ist es
über $16$ Versuche äußerst unwahrscheinlich, höchstens eine grüne Kugel zu
ziehen. Dies zeigt sich auch am \emph{Erwartungswert}
(Abschnitt~\ref{sec:erwartungswert}): Für eine binomialverteilte Zufallsvariable
$X = \sum_{i=1}^{n} X_i$ als Summe von $n$ unabhängigen Bernoulli-Variablen mit
$\mathrm{E}\{X_i\} = P_g$ folgt aus der Linearität des Erwartungswerts

$$
\mu_X = \mathrm{E}\{X\} = \sum_{k=0}^{n} k\, p_X(k) = n\,P_g
$$

und damit konkret

$$
\mu_X = n\,P_g = 16 \cdot 0{,}4 = 6{,}4 \text{ grüne Kugeln}.
$$

Der Erwartungswert liegt also weit über $1$, weshalb $P(X \le 1)$ verschwindend
klein ausfällt.

\end{beispiel}

% ================================================================
\subsection*{Aufgaben 8.1 – 8.8}
% ================================================================

Die folgenden Aufgaben vertiefen den Satz von Bayes, den Satz von der totalen
Wahrscheinlichkeit sowie die Binomial- und die hypergeometrische Verteilung.

% ----------------------------------------------------------------
\subsubsection*{8.1 – Fahren ohne gültigen Fahrschein}

\begin{beispiel}[Satz von Bayes -- Schwarzfahrer]

\textbf{Gegeben:} In Bochum wird zu $7\,\%$ schwarzgefahren, also
$P(S) = 0{,}07$ (Schwarzfahrer) und $P(E) = 0{,}93$ (ehrlicher Fahrgast).
Von den Schwarzfahrern haben $80\,\%$ gar keine Fahrkarte, die übrigen
$20\,\%$ besitzen eine gefälschte bzw.\ illegal besorgte Fahrkarte. Von den
ehrlichen Fahrgästen haben $5\,\%$ ihren Fahrschein vergessen. Betrachtet wird
das Ereignis $K$: „der kontrollierte Fahrgast kann keinen Fahrschein
vorzeigen``. Ein Schwarzfahrer mit gefälschtem Schein kann diesen vorzeigen,
sodass

$$
P(K \mid S) = 0{,}80, \qquad P(K \mid E) = 0{,}05 .
$$

\textbf{Gesucht:} a) $P(S \mid K)$ und b) $P(E \mid K)$.

\textbf{Formel} (Satz von Bayes mit totaler Wahrscheinlichkeit im Nenner,
Abschnitte~\ref{sec:satz-von-bayes} und~\ref{sec:totale-wahrscheinlichkeit}):

$$
P(S \mid K) = \frac{P(K \mid S)\, P(S)}{P(K \mid S)\, P(S) + P(K \mid E)\, P(E)}
$$

\textbf{Rechnung:}

Zunächst die totale Wahrscheinlichkeit für „kein Fahrschein vorzeigbar``:

$$
P(K) = 0{,}80 \cdot 0{,}07 + 0{,}05 \cdot 0{,}93
     = 0{,}056 + 0{,}0465 = 0{,}1025
$$

a) Rückschluss auf den Schwarzfahrer:

$$
P(S \mid K) = \frac{0{,}056}{0{,}1025} = \frac{112}{205} \approx 0{,}5463
$$

b) Rückschluss auf den ehrlichen Fahrgast:

$$
P(E \mid K) = \frac{0{,}0465}{0{,}1025} = \frac{93}{205} \approx 0{,}4537
$$

Probe: $\tfrac{112}{205} + \tfrac{93}{205} = \tfrac{205}{205} = 1$ \checkmark

\textbf{Lösung:}

$$
\boxed{P(S \mid K) = \frac{112}{205} \approx 0{,}5463, \qquad
P(E \mid K) = \frac{93}{205} \approx 0{,}4537}
$$

Wer keinen Fahrschein vorzeigen kann, ist mit knapp $55\,\%$ ein Schwarzfahrer,
obwohl der Anteil der Schwarzfahrer insgesamt nur $7\,\%$ beträgt.

\end{beispiel}

% ----------------------------------------------------------------
\subsubsection*{8.2 – Bestehen einer Klausur}

\begin{beispiel}[Totale Wahrscheinlichkeit und Bayes -- Klausur]

\textbf{Gegeben:} $V$ = „Student ist vorbereitet`` mit $P(V) = 0{,}25$,
$P(\bar V) = 0{,}75$. $B$ = „Prüfung bestanden`` mit den bedingten
Wahrscheinlichkeiten

$$
P(B \mid V) = 0{,}98, \qquad P(B \mid \bar V) = 0{,}01 .
$$

\textbf{Gesucht:} a) $P(B)$, b) $P(\bar B \cap V)$, c) $P(V \mid \bar B)$.

\textbf{Formel} (Satz von der totalen Wahrscheinlichkeit
\ref{sec:totale-wahrscheinlichkeit}, Multiplikationssatz und Bayes
\ref{sec:satz-von-bayes}):

$$
P(B) = P(B \mid V)\,P(V) + P(B \mid \bar V)\,P(\bar V), \qquad
P(V \mid \bar B) = \frac{P(\bar B \cap V)}{P(\bar B)}
$$

\textbf{Rechnung:}

a)
$$
P(B) = 0{,}98 \cdot 0{,}25 + 0{,}01 \cdot 0{,}75 = 0{,}245 + 0{,}0075 = 0{,}2525
$$

b) Ein vorbereiteter Student fällt mit $P(\bar B \mid V) = 0{,}02$ durch:

$$
P(\bar B \cap V) = P(\bar B \mid V)\, P(V) = 0{,}02 \cdot 0{,}25 = 0{,}005
$$

c) Mit $P(\bar B) = 1 - P(B) = 0{,}7475$:

$$
P(V \mid \bar B) = \frac{0{,}005}{0{,}7475} = \frac{2}{299} \approx 0{,}00669
$$

\textbf{Lösung:}

$$
\boxed{P(B) = 0{,}2525, \qquad P(\bar B \cap V) = 0{,}005, \qquad
P(V \mid \bar B) = \frac{2}{299} \approx 0{,}00669}
$$

Wer durchfällt, war nur mit rund $0{,}67\,\%$ Wahrscheinlichkeit vorbereitet --
Durchfallen spricht sehr stark für mangelnde Vorbereitung.

\end{beispiel}

% ----------------------------------------------------------------
\subsubsection*{8.3 – Bitfehlerwahrscheinlichkeiten}

\begin{beispiel}[Binomialverteilung -- Übertragungsfehler]

\textbf{Gegeben:} Ein Übertragungskanal überträgt jedes Bit unabhängig mit der
Bitfehlerwahrscheinlichkeit $P_e$. Ein Datenpaket besteht aus $n$ Bits. Die
Anzahl $X$ der fehlerhaften Bits ist damit binomialverteilt (Bernoulli-Kette).

\textbf{Gesucht:} Wahrscheinlichkeit für a) keinen Fehler, b) genau einen
Fehler, c) genau $k$ Fehler, d) höchstens $3$ Fehler.

\textbf{Formel} (Binomialverteilung, Abschnitt~\ref{sec:binomialverteilung}):

$$
P(X = k) = \binom{n}{k}\, P_e^{\,k}\,(1 - P_e)^{\,n-k}
$$

\textbf{Rechnung:}

a) Kein Fehler ($k = 0$):

$$
P(X = 0) = \binom{n}{0}\, P_e^{0}\,(1 - P_e)^{n} = (1 - P_e)^{n}
$$

b) Genau ein Fehler ($k = 1$):

$$
P(X = 1) = \binom{n}{1}\, P_e^{1}\,(1 - P_e)^{n-1} = n\,P_e\,(1 - P_e)^{n-1}
$$

c) Genau $k$ Fehler:

$$
P(X = k) = \binom{n}{k}\, P_e^{\,k}\,(1 - P_e)^{\,n-k}
$$

d) Höchstens $3$ Fehler ($X \le 3$):

$$
P(X \le 3) = \sum_{k=0}^{3} \binom{n}{k}\, P_e^{\,k}\,(1 - P_e)^{\,n-k}
$$

\textbf{Lösung:}

$$
\boxed{%
\begin{aligned}
P(X = 0) &= (1 - P_e)^{n}, \\
P(X = 1) &= n\,P_e\,(1 - P_e)^{n-1}, \\
P(X = k) &= \binom{n}{k}\, P_e^{\,k}\,(1 - P_e)^{\,n-k}, \\
P(X \le 3) &= \sum_{k=0}^{3} \binom{n}{k}\, P_e^{\,k}\,(1 - P_e)^{\,n-k} .
\end{aligned}}
$$

\end{beispiel}

% ----------------------------------------------------------------
\subsubsection*{8.4 – Restposten ICs}

\begin{beispiel}[Hypergeometrische Verteilung -- defekte ICs]

\textbf{Gegeben:} Ein Posten enthält $N = 10$ ICs, davon $D = 4$ defekt und
$6$ intakt. Es werden \emph{ohne Zurücklegen} $n = 4$ ICs zufällig ausgewählt.
Die Anzahl $X$ der defekten unter den gezogenen ICs ist hypergeometrisch
verteilt.

\textbf{Gesucht:} $P(X = 4)$, $P(X = 1)$, $P(X = 2)$.

\textbf{Formel} (Laplace-Wahrscheinlichkeit über Kombinationen,
Abschnitt~\ref{sec:kombinatorik}):

$$
P(X = m) = \frac{\dbinom{D}{m}\dbinom{N-D}{\,n-m\,}}{\dbinom{N}{n}}
$$

\textbf{Rechnung:}

Anzahl aller Auswahlen: $\dbinom{10}{4} = 210$.

a) Alle vier defekt ($m = 4$):

$$
P(X = 4) = \frac{\dbinom{4}{4}\dbinom{6}{0}}{\dbinom{10}{4}}
         = \frac{1 \cdot 1}{210} = \frac{1}{210} \approx 0{,}00476
$$

b) Genau eins defekt ($m = 1$):

$$
P(X = 1) = \frac{\dbinom{4}{1}\dbinom{6}{3}}{\dbinom{10}{4}}
         = \frac{4 \cdot 20}{210} = \frac{80}{210} = \frac{8}{21} \approx 0{,}3810
$$

c) Genau zwei defekt ($m = 2$):

$$
P(X = 2) = \frac{\dbinom{4}{2}\dbinom{6}{2}}{\dbinom{10}{4}}
         = \frac{6 \cdot 15}{210} = \frac{90}{210} = \frac{3}{7} \approx 0{,}4286
$$

\textbf{Lösung:}

$$
\boxed{P(X = 4) = \frac{1}{210} \approx 0{,}00476, \quad
P(X = 1) = \frac{8}{21} \approx 0{,}3810, \quad
P(X = 2) = \frac{3}{7} \approx 0{,}4286}
$$

\end{beispiel}

% ----------------------------------------------------------------
\subsubsection*{8.5 – Radarstation}

\begin{beispiel}[Satz von Bayes -- Radarstation]

\textbf{Gegeben:} $N$ = „störungsüberlagertes Nutzsignal`` mit $P(N) = 0{,}90$,
$T$ = „reine Störung`` mit $P(T) = 0{,}10$. Die Anzeige $A$ = „Nutzsignal
angezeigt`` tritt ein mit

$$
P(A \mid N) = 0{,}98, \qquad P(A \mid T) = 0{,}10 .
$$

\textbf{Gesucht:} $P(N \mid A)$ -- Wahrscheinlichkeit, dass ein als Nutzsignal
angezeigtes Signal wirklich ein Nutzsignal ist.

\textbf{Formel} (Satz von Bayes, Abschnitt~\ref{sec:satz-von-bayes}):

$$
P(N \mid A) = \frac{P(A \mid N)\,P(N)}{P(A \mid N)\,P(N) + P(A \mid T)\,P(T)}
$$

\textbf{Rechnung:}

$$
P(A) = 0{,}98 \cdot 0{,}90 + 0{,}10 \cdot 0{,}10 = 0{,}882 + 0{,}010 = 0{,}892
$$

$$
P(N \mid A) = \frac{0{,}882}{0{,}892} = \frac{441}{446} \approx 0{,}9888
$$

\textbf{Lösung:}

$$
\boxed{P(N \mid A) = \frac{441}{446} \approx 0{,}9888}
$$

Eine „Nutzsignal``-Anzeige ist mit rund $98{,}9\,\%$ zuverlässig, da echte
Nutzsignale sowohl häufig als auch fast sicher korrekt angezeigt werden.

\end{beispiel}

% ----------------------------------------------------------------
\subsubsection*{8.6 – Kontrolle am Flughafen}

\begin{beispiel}[Satz von Bayes -- Sicherheitskontrolle]

\textbf{Gegeben:} $V$ = „verbotene Gegenstände im Handgepäck`` mit
$P(V) = 0{,}20$, $P(\bar V) = 0{,}80$. Das Scan-Gerät löst Alarm $A$ aus mit

$$
P(A \mid V) = 0{,}98, \qquad P(A \mid \bar V) = 0{,}01 \quad (\text{Fehlalarm}).
$$

\textbf{Gesucht:} a) $P(A)$, b) $P(V \mid A)$, c) Anzahl der Fehlalarme bei
$1\,000\,000$ Fluggästen, d) $P(\bar V \mid A)$ (ohne Ergebnis aus b).

\textbf{Formel} (totale Wahrscheinlichkeit \ref{sec:totale-wahrscheinlichkeit}
und Bayes \ref{sec:satz-von-bayes}):

$$
P(A) = P(A \mid V)\,P(V) + P(A \mid \bar V)\,P(\bar V), \qquad
P(V \mid A) = \frac{P(A \mid V)\,P(V)}{P(A)}
$$

\textbf{Rechnung:}

a)
$$
P(A) = 0{,}98 \cdot 0{,}20 + 0{,}01 \cdot 0{,}80 = 0{,}196 + 0{,}008 = 0{,}204
$$

b)
$$
P(V \mid A) = \frac{0{,}196}{0{,}204} = \frac{49}{51} \approx 0{,}9608
$$

c) Ein Fehlalarm ist das Ereignis $A \cap \bar V$ mit

$$
P(A \cap \bar V) = P(A \mid \bar V)\,P(\bar V) = 0{,}01 \cdot 0{,}80 = 0{,}008 .
$$

Bei $1\,000\,000$ Fluggästen also

$$
0{,}008 \cdot 10^{6} = 8\,000 \text{ Fehlalarme}.
$$

d) Direkt über Bayes (unabhängig von b):

$$
P(\bar V \mid A) = \frac{P(A \mid \bar V)\,P(\bar V)}{P(A)}
                 = \frac{0{,}008}{0{,}204} = \frac{2}{51} \approx 0{,}0392
$$

Probe: $\tfrac{49}{51} + \tfrac{2}{51} = 1$ \checkmark

\textbf{Lösung:}

$$
\boxed{P(A) = 0{,}204, \quad P(V \mid A) = \frac{49}{51} \approx 0{,}9608, \quad
8\,000 \text{ Fehlalarme}, \quad P(\bar V \mid A) = \frac{2}{51} \approx 0{,}0392}
$$

\end{beispiel}

% ----------------------------------------------------------------
\subsubsection*{8.7 – Diabetes-Test}

\begin{beispiel}[Satz von Bayes -- Diabetes-Test]

\textbf{Gegeben:} $K$ = „an Diabetes erkrankt`` mit $P(K) = 0{,}072$,
$P(\bar K) = 0{,}928$. Der Test $T$ (positiv) erfüllt

$$
P(T \mid K) = 0{,}95, \qquad P(\bar K \cap T) = 0{,}10 .
$$

\textbf{Gesucht:} $P(K \mid T)$ -- Wahrscheinlichkeit, tatsächlich krank zu
sein, wenn der Test positiv ausfällt.

\textbf{Formel} (Satz von Bayes, Abschnitt~\ref{sec:satz-von-bayes}):

$$
P(K \mid T) = \frac{P(T \cap K)}{P(T)}, \qquad
P(T) = P(T \mid K)\,P(K) + P(\bar K \cap T)
$$

\textbf{Rechnung:}

Wahrscheinlichkeit für richtig-positiv:

$$
P(T \cap K) = P(T \mid K)\,P(K) = 0{,}95 \cdot 0{,}072 = 0{,}0684
$$

Totale Wahrscheinlichkeit eines positiven Tests (der falsch-positive Anteil ist
mit $0{,}10$ direkt gegeben):

$$
P(T) = 0{,}0684 + 0{,}10 = 0{,}1684
$$

$$
P(K \mid T) = \frac{0{,}0684}{0{,}1684} = \frac{171}{421} \approx 0{,}4062
$$

\textbf{Lösung:}

$$
\boxed{P(K \mid T) = \frac{171}{421} \approx 0{,}4062}
$$

Trotz eines positiven Tests ist man nur mit rund $40{,}6\,\%$ tatsächlich krank,
da die große gesunde Bevölkerungsgruppe viele falsch-positive Ergebnisse liefert
(Basisraten-Effekt).

\end{beispiel}

% ----------------------------------------------------------------
\subsubsection*{8.8 – Relative Häufigkeiten}

\begin{beispiel}[Relative Häufigkeiten im Schwarz-Weiß-Bild]

\textbf{Gegeben:} Ein Schwarz-Weiß-Bild (Bild~1) besteht ausschließlich aus
schwarzen und weißen Bildpunkten. Aus früheren Untersuchungen ist die
Wahrscheinlichkeit für einen schwarzen Bildpunkt mit $P(S) = 0{,}3$ bekannt.

\emph{Annahme (aus dem gerenderten Bild abgelesen):} Das Punktraster ist
quadratisch mit $20$ Spalten und $20$ Zeilen. Ein exaktes pixelgenaues
Auszählen ist aus der Vorlage nicht zuverlässig möglich; die im Folgenden mit
\glqq{}angenommen\grqq{} gekennzeichneten Zahlen sind daher plausible,
konsistent gewählte Werte, mit denen das methodische Vorgehen demonstriert
wird.

\textbf{Gesucht:} a) Anzahl aller Bildpunkte, b) relative Häufigkeiten
schwarzer und weißer Bildpunkte, c) Prüfung der Binomial-Näherung für die
Anzahl schwarzer Punkte pro Zeile, d) akkumulierte relative Häufigkeit über die
Bildzeilen.

\textbf{Formel} (relative Häufigkeit \ref{sec:wahrscheinlichkeit-einfuehrung};
Erwartungswert und Varianz der Binomialverteilung
\ref{sec:binomialverteilung},~\ref{sec:erwartungswert}):

$$
h(A) = \frac{n_A}{N}, \qquad \mu = n\,p, \qquad \sigma^2 = n\,p\,(1-p)
$$

\textbf{Rechnung:}

\emph{a) Anzahl aller Bildpunkte.} Bei $20$ Spalten und $20$ Zeilen:

$$
N = 20 \cdot 20 = 400 \text{ Bildpunkte (Annahme quadratisches Raster).}
$$

\emph{b) Relative Häufigkeiten.} Aus $P(S) = 0{,}3$ folgt für die erwartete
Anzahl schwarzer Punkte $n_S \approx P(S)\cdot N = 0{,}3 \cdot 400 = 120$ und
somit

$$
h(S) = \frac{120}{400} = 0{,}3, \qquad
h(W) = \frac{280}{400} = 0{,}7 = 1 - h(S).
$$

\emph{c) Binomial-Näherung pro Zeile.} Eine Bildzeile besteht aus $n = 20$
Punkten; deutet man jeden Punkt als unabhängiges Bernoulli-Experiment mit
$p = P(S) = 0{,}3$, so ist die Anzahl schwarzer Punkte pro Zeile
binomialverteilt mit den theoretischen Kennwerten

$$
\mu = n\,p = 20 \cdot 0{,}3 = 6, \qquad
\sigma^2 = n\,p\,(1-p) = 20 \cdot 0{,}3 \cdot 0{,}7 = 4{,}2 .
$$

Für die $20$ Zeilen werden die (angenommenen) schwarzen Punktzahlen

$$
x_i = 8,\,4,\,9,\,3,\,7,\,5,\,10,\,2,\,6,\,7,\,5,\,6,\,8,\,4,\,6,\,5,\,7,\,3,\,6,\,9
$$

herangezogen (Summe $\sum x_i = 120$). Damit ergibt sich der empirische
Mittelwert

$$
\bar{x} = \frac{1}{20}\sum_{i=1}^{20} x_i = \frac{120}{20} = 6
$$

und die empirische Varianz

$$
s^2 = \frac{1}{20}\sum_{i=1}^{20} (x_i - \bar{x})^2 = \frac{90}{20} = 4{,}5 .
$$

Der Vergleich $\bar{x} = 6 = \mu$ und $s^2 = 4{,}5 \approx 4{,}2 = \sigma^2$ zeigt
eine sehr gute Übereinstimmung. Die Anzahl schwarzer Bildpunkte pro Zeile lässt
sich damit gut durch eine Binomialverteilung $\mathcal{B}(n = 20,\,p = 0{,}3)$
annähern.

\emph{d) Akkumulierte relative Häufigkeit.} Werden die schwarzen Bildpunkte von
oben nach unten aufsummiert und durch die bis dahin betrachtete Punktzahl
($\text{Zeile} \cdot 20$) geteilt, so ergibt sich die kumulierte relative
Häufigkeit $h_m = \tfrac{1}{20\,m}\sum_{i=1}^{m} x_i$. Diese pendelt sich mit
wachsender Zeilenzahl auf den Wert $P(S) = 0{,}3$ ein (empirisches Gesetz der
großen Zahlen):

\begin{center}
\begin{tikzpicture}
\begin{axis}[
  width=0.9\linewidth, height=6cm,
  xlabel={Bildzeile $m$}, ylabel={akkum.\ rel.\ Häufigkeit $h_m$},
  xmin=0, xmax=21, ymin=0.25, ymax=0.42,
  xtick={2,4,6,8,10,12,14,16,18,20},
  ytick={0.28,0.30,0.32,0.34,0.36,0.38,0.40},
  grid=both, grid style={gray!25},
  legend pos=north east, legend cell align=left,
]
% angenommene kumulierte relative Häufigkeit
\addplot[thick, mark=*, mark size=1.2pt, color=accent] coordinates {
  (1,0.400) (2,0.300) (3,0.350) (4,0.300) (5,0.310)
  (6,0.300) (7,0.329) (8,0.300) (9,0.300) (10,0.305)
  (11,0.300) (12,0.300) (13,0.308) (14,0.300) (15,0.300)
  (16,0.297) (17,0.300) (18,0.292) (19,0.292) (20,0.300)
};
\addlegendentry{$h_m$ (angenommen)}
\addplot[dashed, thick, color=satzred-fr, domain=0:21] {0.3};
\addlegendentry{$P(S) = 0{,}3$}
\end{axis}
\end{tikzpicture}
\end{center}

\textbf{Lösung:}

$$
\boxed{N = 400, \quad h(S) = 0{,}3, \quad h(W) = 0{,}7, \quad
\bar{x} = 6 \approx \mu, \quad s^2 = 4{,}5 \approx \sigma^2 = 4{,}2}
$$

Die Anzahl schwarzer Bildpunkte pro Zeile ist gut binomialverteilt, und die
akkumulierte relative Häufigkeit konvergiert gegen $P(S) = 0{,}3$.

\end{beispiel}
